% BEGIN PROFESSOR FEEDBACK PF35: Foundations - Graph-theory citations
\section{Graph-theory foundations}
\label{sec:graph}

A graph is a mathematical structure consisting of a set of nodes $V$ and a set 
of edges $E$, where nodes represent entities and edges represent pairwise 
relations between them \citep{newman2010networks}. In the context of this
thesis, nodes correspond to Ethereum addresses and edges correspond to 
transactions between them.

Because transactions have a clearly defined sender and receiver, the 
appropriate structure is a \emph{directed graph}, where an edge $(u, v) \in E$ 
represents a directed relation from node $u$ to node $v$ 
\citep{newman2010networks}. Since the same pair of addresses may transact
multiple times, each with a distinct value and timestamp, the data is more 
precisely represented as a \emph{directed multigraph}, in which multiple edges 
between the same ordered pair of nodes are permitted 
\citep{paranjape2017motifs, sanjalawe2023atd}.

For a node $v$, the \emph{in-degree} $d^-(v)$ counts the number of incoming 
edges and the \emph{out-degree} $d^+(v)$ counts the number of outgoing edges 
\citep{newman2010networks}. These degree measures serve as basic node-level
features and are used as baselines in the classification experiments in 
Chapter~\ref{ch:method}.

The directed multigraph provides the structural foundation for the graph model 
used throughout this thesis. The following section extends this static 
representation by incorporating the timestamps of transactions, which is 
necessary to capture the temporal ordering of interactions.
% END PROFESSOR FEEDBACK PF35
